\documentclass[11pt]{article}
\usepackage[margin=2.6cm]{geometry}
\usepackage{amsmath,amssymb,amsthm}
\usepackage[hidelinks]{hyperref}
\newtheorem{theorem}{Theorem}
\newtheorem{proposition}[theorem]{Proposition}
\newtheorem{lemma}[theorem]{Lemma}
\newtheorem{corollary}[theorem]{Corollary}
\theoremstyle{remark}
\newtheorem{remark}[theorem]{Remark}
\newcommand{\C}{\mathbb{C}}
\newcommand{\Hess}{\operatorname{Hess}}
\newcommand{\adj}{\operatorname{adj}}

\title{The four-dimensional Hessian conjecture for quintic potentials\\ with a ternary cone top}
\author{David Niedbala Giraudin\thanks{ORCID 0009-0009-1526-1178. Computations and writing assisted by Claude (Anthropic).}}
\date{5 November 2026}

\begin{document}
\maketitle

\begin{abstract}
The Hessian conjecture $\mathrm{HC}_n$ asserts that if $f\in\C[x_1,\dots,x_n]$ has $\det\Hess(f)\in\C^\times$, then $\nabla f$ is a polynomial automorphism. It is known for $n\le 3$, false for $n\ge 5$, and open for $n=4$, where it implies the plane Jacobian conjecture. Z.~Ni proved $\mathrm{HC}_4$ for potentials of degree at most four. We prove $\mathrm{HC}_4$ for every potential of degree five whose top homogeneous form depends on exactly three linear forms. The proof combines a weighted analysis of the Hessian determinant with exact computer-algebra certificates (Gr\"obner bases and Rabinowitsch tests over $\mathbb{Q}$); one branch reduces to plane Jacobian pairs of degrees $(4,3)$, which are excluded by Magnus' theorem. A self-contained verifier is supplied as an ancillary file. The binary and unary cone tops in degree five remain open.
\end{abstract}

\section{Introduction}
For $f\in\C[x]=\C[x_1,\dots,x_n]$ write $\Hess(f)=(\partial_i\partial_j f)$. The Hessian conjecture $\mathrm{HC}_n$ \cite{Meng06} states: if $\det\Hess(f)\in\C^\times$ then $\nabla f\colon\C^n\to\C^n$ is a polynomial automorphism. It holds for $n=2$ \cite{Dillen} and $n=3$ \cite{deBondt3}; it fails for every $n\ge5$ \cite{MengYang}; for $n=4$ it is open and implies the plane Jacobian conjecture by the doubling construction $\Phi(x,y)=y^{\mathsf T}F(x)$, for which $\det\Hess\Phi=(\det JF)^2$ \cite{MengYang,Ni}. Ni \cite{Ni} proved $\mathrm{HC}_4$ for $\deg f\le 4$. In degree five the determinant has more homogeneous layers and the arguments of \cite{Ni} do not transfer directly (\cite[\S8]{Ni}); partial structural reductions for a degree-five branch are recorded in \cite{EP}.

Let $f\in\C[x_1,\dots,x_4]$ have degree $5$ and $\det\Hess(f)=c\in\C^\times$. Its top form $f_5$ has $\det\Hess(f_5)=0$, hence by Gordan--Noether \cite{GN,deBondtQT} depends on at most three linear forms. We say the top is a \emph{ternary cone} if, after a linear change of coordinates, $f_5=\varphi(x_1,x_2,x_3)$ with $\det\Hess_3(\varphi)\not\equiv0$ (equivalently, $f_5$ depends on exactly three linear forms).

\begin{theorem}\label{thm:main}
Let $f\in\C[x_1,x_2,x_3,x_4]$ have degree five, $\det\Hess(f)\in\C^\times$, and a ternary cone top. Then $\nabla f$ is a polynomial automorphism of $\C^4$.
\end{theorem}

\begin{remark}
The class of Theorem~\ref{thm:main} contains doubled potentials $F(x_1,x_2)+x_3a(x_1,x_2)+x_4b(x_1,x_2)$, whose Hessian determinant is $J(a,b)^2$ (Section~\ref{sec:D}). Hence no proof of the quintic case can avoid plane Jacobian pairs of low degree; here only pairs of degrees $(4,3)$ occur, and Magnus' theorem excludes them.
\end{remark}

\paragraph{Status.} The statements labelled \emph{[C]} below are certified by exact computation (identities of polynomials over $\mathbb{Q}$ with symbolic coefficients, Gr\"obner bases, and the Rabinowitsch criterion $1\in I+(1-Tg)\Leftrightarrow g\in\sqrt I$). Each such statement corresponds to a check of the ancillary verifier \texttt{verify\_hc4\_ternary.py} (Section~\ref{sec:cert}); all 29 checks pass. The paper has not been refereed.

\section{Tools}
Throughout, $x'=(x_1,x_2,x_3)$, $s=x_4$, and $A=\Hess_3(\varphi)$. Linear changes of coordinates preserve both hypothesis and conclusion, since $\Hess(f\circ T)=T^{\mathsf T}(\Hess f)\circ T\,T$.

\paragraph{Imported results.}
(GN) A form in $n\le4$ variables with vanishing Hessian determinant depends on at most $n-1$ linear forms \cite{GN}, \cite[Thm.~3.1]{deBondtQT}.
(HC$_2$, HC$_3$) \cite{Dillen,deBondt3}, together with transfer to arbitrary fields of characteristic zero and descent of polynomial inverses \cite[Lemmas~2.6, 2.7]{Ni}.
(dBvdE) Over a field of characteristic zero, a bivariate polynomial without terms of degree $\le1$ and with vanishing Hessian determinant is a polynomial in one linear form \cite[Thm.~3.1]{dBvdE}.
(Ni) \cite[Prop.~6.1]{Ni}: if $f=P(x')+sQ(x')+as^2$ with $\deg Q\le2$ and $\det\Hess f\in\C^\times$, then $\nabla f$ is an automorphism.
(Mag) \cite{Magnus}: a plane Jacobian pair whose degrees are coprime is an automorphism.
(JvdK) For a plane automorphism $(F,G)$, $\deg F\mid\deg G$ or $\deg G\mid\deg F$ \cite[Ch.~5]{vdE}.

\begin{lemma}[Schur branch, any $Q$]\label{lem:A}
Let $f=P(x')+sQ(x')+as^2$ with $a\ne0$ and $P,Q$ arbitrary. If $\det\Hess f\in\C^\times$, then $\nabla f$ is a polynomial automorphism.
\end{lemma}
\begin{proof}
With $z=s+Q/(2a)$ and $S_z=P-Q^2/(4a)+zQ$ one has the polynomial identity $\det\Hess f=2a\det\Hess_{x'}(S_z)$ at $z=s+Q/(2a)$ \emph{[C, S6a]}, and $\nabla_{x'}f=\nabla_{x'}S_z$, $\partial_sf=2az$. Thus $\det\Hess_{x'}S_z$ is a nonzero constant; $\mathrm{HC}_3$ over $\C(z)$ and descent to $\C[z]$ give a polynomial inverse $H(z,Y)$ of $\nabla_{x'}S_z$, and $z=Y_4/(2a)$, $x'=H(z,Y_1,Y_2,Y_3)$, $s=z-Q(x')/(2a)$ invert $\nabla f$. This is the argument of \cite[Prop.~6.1, branch $a\ne0$]{Ni}, which does not use $\deg Q\le2$.
\end{proof}

\begin{lemma}[constant direction]\label{lem:H}
Let $F\in\C[x_1,x_2,x_3]$ be homogeneous of degree $d$, nonzero, without terms of degree $\le1$ in $(x_2,x_3)$, and with $\det\Hess_{(x_2,x_3)}F=0$. Then there is a constant linear form $l=\lambda x_2+\mu x_3$ with $F\in\C[x_1][l]$.
\end{lemma}
\begin{proof}
Over $K=\C(x_1)$, (dBvdE) gives $F=G(l)$ with $l=\lambda x_2+\mu x_3$, $\lambda,\mu\in K$, $G\in K[T]$, $\deg G\ge2$. Then $\Hess_{(x_2,x_3)}F=G''(l)\,(\lambda,\mu)^{\mathsf T}(\lambda,\mu)$ has rank exactly one, so its kernel line is determined by $F$. Since $\Hess_{(x_2,x_3)}F(tx)=t^{d-2}\Hess_{(x_2,x_3)}F(x)$, the kernel line at $tx_1$ equals the one at $x_1$; thus the ratio $\lambda/\mu$ (or $\mu/\lambda$), a rational function of $x_1$, is invariant under all dilations and hence constant. After rescaling $\lambda,\mu\in\C$, $l$ is a coordinate and the coefficients of $F$ as a polynomial in $l$ lie in $\C[x_1]$.
\end{proof}

\section{The weighted top}
\begin{proposition}\label{prop:top}
Let $f$ be as in Theorem~\ref{thm:main}, $f=f_5+\dots+f_0$, $f_5=\varphi(x')$. Then:
\begin{enumerate}
\item $[\det\Hess f]_{11}=\det A\cdot\partial_s^2f_4$ \emph{[C, S1b]}; hence $f_4=g_4(x')+s\,h(x')$.
\item With this, $[\det\Hess f]_{10}=\det A\cdot\partial_s^2f_3-(\nabla h)^{\mathsf T}\adj(A)\nabla h$ \emph{[C, S1c]}; writing $f_3=g_3+sh_2+s^2m(x')+es^3$, it follows that $e=0$.
\item For the weights $(1,1,1,2)$, $f$ has no monomial of weight $>5$, its weight-$5$ part is $\psi=\varphi+sh+s^2m$, and the weight-$10$ part of $\det\Hess f$ equals $\det\Hess\psi$ \emph{[C, S2]}. Hence $\det\Hess\psi=0$.
\end{enumerate}
\end{proposition}
\begin{proof}
(1) and (2) are identities with fully symbolic $f_4,f_3,f_2$; as $\det\Hess f$ is constant and $\det A\ne0$, $\partial_s^2f_4=0$, and in (2) the term $6e s\det A$ is the only $s$-dependent one, so $e=0$. (3) After (1) and (2), $f$ contains no $s^2x'^2$, $s^3x'$, $s^4$ (degree 4) and no $s^3$ (degree 3); every other monomial outside $\varphi$ has weight $\le4$, except $s\,x'^3$ (weight 5, in $sh$) and $s^2x'$ (weight 5, in $s^2m$); the identity holds for arbitrary lower terms, and the constant has weight $0<10$.
\end{proof}

\section{Classification of the weighted top}
\begin{proposition}\label{prop:class}
Let $\psi=\varphi(x')+sh(x')+s^2m(x')$ with $\varphi$ a quintic, $h$ a cubic, $m$ linear, $\det\Hess_3\varphi\not\equiv0$ and $\det\Hess\psi=0$. After a linear change of coordinates preserving the shape of $\psi$, one of the following holds:
\begin{itemize}
\item[(N)] $h=m=0$;
\item[(D)] $m=0$ and $\psi=Q(x_1,x_2)+x_1^4x_3+x_1^3s$ with $\partial_2^2Q\ne0$;
\item[(F)] $\psi=x_1(s+x_1x_2)^2+x_1^4x_3$.
\end{itemize}
\end{proposition}
\begin{proof}
\emph{Case $m=0$.} The coefficient of $s^2$ in $\det\Hess(\varphi+sh)$ is $-\tfrac32\,h\det\Hess_3h$, and those of $s^3,s^4$ vanish \emph{[C, S3a, S3b]}. So $h=0$ (case (N)) or, by (GN), $h$ is a binary cubic in two linear forms; up to $\mathrm{GL}_3$, $h\in\{x_1^3,\ x_1^2x_2,\ x_1^3+x_2^3\}$. For $h=x_1^2x_2$ and $h=x_1^3+x_2^3$, every coefficient of $\det A$ lies in the radical of the ideal of coefficients of $\det\Hess\psi$ \emph{[C, S3c]}: excluded. For $h=x_1^3$, $\det\Hess\psi=-9x_1^4\det\Hess_{(x_2,x_3)}\varphi$ \emph{[C, S3d]}. Apply Lemma~\ref{lem:H} to the part of $\varphi$ of degree $\ge2$ in $(x_2,x_3)$, which is nonzero (otherwise $A$ has a zero $2\times2$ block and $\det A=0$). After a rotation of $(x_2,x_3)$, $\varphi=Q(x_1,x_2)+x_1^4(\beta_2x_2+\alpha x_3)+\beta_0x_1^5$, and $\det A=-16\alpha^2x_1^6\,\partial_2^2Q$ \emph{[C, S8c]}, so $\alpha\ne0\ne\partial_2^2Q$. Replacing $x_3$ by $\alpha x_3+\beta_2x_2+\beta_0x_1$ gives (D).

\emph{Case $m\ne0$.} Take $m=x_1$. The coefficient of $s^4$ is $-4\det\Hess_{(x_2,x_3)}h$ \emph{[C, S4a]}; Lemma~\ref{lem:H} (applied to the part of $h$ of degree $\ge2$ in $(x_2,x_3)$, if nonzero) and a rotation of $(x_2,x_3)$ give $h=a x_2^3+bx_1x_2^2+x_1^2(cx_2+c'x_3)+ex_1^3$. If $a\ne0$, rescale to $a=1$; if $a=0\ne b$, rescale to $b=1$ and shift $x_2\mapsto x_2-\tfrac c2x_1$ to get $c=0$; in both cases $\det A$ vanishes on all solutions \emph{[C, S4 (I), (II)]}. If $a=b=0$ and $(c,c')\ne0$, a rotation, rescaling and shift give $h=x_1^2x_2$; otherwise $h=ex_1^3$, which is excluded \emph{[C, S4 (III.2)]}. For $h=x_1^2x_2$ the solutions with $\det A\not\equiv0$ are exactly $\varphi=\tfrac14x_1^3x_2^2+p_{18}x_1^4x_3+p_{19}x_1^4x_2+p_{20}x_1^5$ \emph{[C, S4 (III.1), S8a]}, with $\det A=-8p_{18}^2x_1^9$ \emph{[C, S8b]}, so $p_{18}\ne0$; the substitution $x_3\leftarrow(x_3-p_{19}x_2-p_{20}x_1)/p_{18}$, $x_2\mapsto2x_2$ gives (F) \emph{[C, S8d]}. All changes used fix $x_1$ and act on $x'$ only, so they preserve $m=x_1$ and the shape of $\psi$.
\end{proof}

\section{Closing the three cases}
\subsection{Case (N)} Here $f$ has no monomials $s\,x'^3$, $s^2x'$, $s^3$, so $f=P(x')+sQ(x')+as^2$ with $\deg Q\le2$, and (Ni) applies.

\subsection{Case (F)}
\begin{proposition}
No quintic $f$ with weight-$5$ part $x_1(s+x_1x_2)^2+x_1^4x_3$ has $\det\Hess f\in\C^\times$.
\end{proposition}
\begin{proof}
Write $f=\psi_F+\sum_\mu c_\mu\mu$ over all monomials $\mu$ of weight $\le4$ and degree $\ge2$ (41 unknowns). The constant term of $\det\Hess f$ lies in the radical of the ideal generated by the coefficients of all non-constant monomials \emph{[C, S5]}.
\end{proof}

\subsection{Case (D)}\label{sec:D}
\begin{proposition}
Let $f$ be quintic with weight-$5$ part $Q(x_1,x_2)+x_1^4x_3+x_1^3s$, $\partial_2^2Q\ne0$, and $\det\Hess f=c\in\C^\times$. Then $\nabla f$ is an automorphism.
\end{proposition}
\begin{proof}
By Proposition~\ref{prop:top} and $m=0$, $f=P(x')+sR(x')+as^2$ with $R=x_1^3+(\text{terms of degree}\le2)$. If $a\ne0$, Lemma~\ref{lem:A} applies. Let $a=0$; then $c=-(\nabla R)^{\mathsf T}\adj(\Hess P+s\Hess R)\nabla R$.

(i) The coefficient of $s^2$ forces $(\partial_2R,\partial_3R)$ to be proportional \emph{[C, S6b]}, so $R=\tilde R(x_1,w)$ for a linear form $w$ in $(x_2,x_3)$. Complete to linear coordinates $(x_1,w,v)$ with $\partial_v x_1=\partial_v w=0$; thus $v=(0,v_2,v_3)$.

(ii) The coefficient of $s$ equals $-P_{vv}\cdot\mathrm{Br}(\tilde R)$, $\mathrm{Br}(R)=R_1^2R_{ww}-2R_1R_wR_{1w}+R_w^2R_{11}$ \emph{[C, S6c]}. As $\C[x]$ is a domain, $\mathrm{Br}=0$ or $P_{vv}=0$.

(iii) If $\mathrm{Br}=0$, then $\tilde R=g(x_1)$ \emph{[C, S6d]} and $c=-g'(x_1)^2\det\Hess_{(w,v)}P$ \emph{[C, S6e]}, impossible since $\deg g'=2$.

(iv) If $P_{vv}=0$: the degree-5 part of $P_{vv}$ is $\partial_v^2(Q+x_1^4x_3)=v_2^2\,\partial_2^2Q$ \emph{[C, S8e]}, so $v_2=0$, $v=e_3$, $R=R(x_1,x_2)$ and $P=P_0(x_1,x_2)+x_3P_1(x_1,x_2)$ with $P_1=x_1^4+(\deg\le3)$. Then $c=J(P_1,R)^2$ \emph{[C, S6f]}, so $(P_1,R)$ is a plane Jacobian pair of degrees $(4,3)$. By (Mag) it is an automorphism, contradicting (JvdK).
\end{proof}

\begin{proof}[Proof of Theorem~\ref{thm:main}]
Propositions~\ref{prop:top} and \ref{prop:class} leave the cases (N), (D), (F), settled in this section.
\end{proof}

\section{Certificates}\label{sec:cert}
The ancillary file \texttt{verify\_hc4\_ternary.py} (Python with \texttt{python-flint}, \texttt{sympy}, and Singular through Sage/passagemath) performs 29 checks, labelled S1a--S8e as cited above, and prints \texttt{ALL CHECKS PASSED}; it runs in a few minutes. Identities are verified with fully symbolic coefficients (or with undetermined functions where stated). Exclusions use the Rabinowitsch criterion over $\mathbb{Q}$; since a Gr\"obner basis equal to $\{1\}$ over $\mathbb{Q}$ is also $\{1\}$ over $\C$, the conclusions hold over $\C$. The classification of S4 (III.1) is certified twice, by minimal associated primes and by Rabinowitsch tests (S8a).

\paragraph{Scope.} Theorem~\ref{thm:main} does not treat quintic potentials whose top form depends on two linear forms or on one; these cases, and hence $\mathrm{HC}_4$ in degree five, remain open.

\begin{thebibliography}{99}
\bibitem{deBondt3} M.~de~Bondt, Polynomials with constant Hessian determinants in dimension three, J.~Pure Appl.\ Algebra 219 (2015), no.~9, 3743--3754.
\bibitem{deBondtQT} M.~de~Bondt, Quasi-translations and singular Hessians, Colloq.\ Math.\ 152 (2018), no.~2, 175--198.
\bibitem{dBvdE} M.~de~Bondt and A.~van~den~Essen, Singular Hessians, J.~Algebra 282 (2004), no.~1, 195--204.
\bibitem{Dillen} F.~Dillen, Polynomials with constant Hessian determinant, J.~Pure Appl.\ Algebra 71 (1991), no.~1, 13--18.
\bibitem{EP} Anonymous, Structural reductions toward the quartic Hessian conjecture in dimension four, v0.1.0-candidate, Zenodo, doi:10.5281/zenodo.22216400 (2026).
\bibitem{vdE} A.~van~den~Essen, Polynomial Automorphisms and the Jacobian Conjecture, Progr.\ Math.\ 190, Birkh\"auser, 2000.
\bibitem{GN} P.~Gordan and M.~N\"other, \"Uber die algebraischen Formen, deren Hesse'sche Determinante identisch verschwindet, Math.\ Ann.\ 10 (1876), 547--568.
\bibitem{Magnus} A.~Magnus, On polynomial solutions of a differential equation, Math.\ Scand.\ 3 (1955), 255--260.
\bibitem{Meng06} G.~Meng, Legendre transform, Hessian conjecture and tree formula, Appl.\ Math.\ Lett.\ 19 (2006), no.~6, 503--510.
\bibitem{MengYang} G.~Meng and L.~Yang, A five-variable counterexample to the Hessian conjecture, and the low-dimensional status of the Jacobian and Hessian conjectures, arXiv:2607.22198 (2026).
\bibitem{Ni} Z.~Ni, The quartic Hessian conjecture in dimension four, arXiv:2608.14217 (2026).
\end{thebibliography}
\end{document}
